% Options for packages loaded elsewhere
\PassOptionsToPackage{unicode}{hyperref}
\PassOptionsToPackage{hyphens}{url}
\documentclass[
]{article}
\usepackage{xcolor}
\usepackage[margin=1in]{geometry}
\usepackage{amsmath,amssymb}
\setcounter{secnumdepth}{-\maxdimen} % remove section numbering
\usepackage{iftex}
\ifPDFTeX
  \usepackage[T1]{fontenc}
  \usepackage[utf8]{inputenc}
  \usepackage{textcomp} % provide euro and other symbols
\else % if luatex or xetex
  \usepackage{unicode-math} % this also loads fontspec
  \defaultfontfeatures{Scale=MatchLowercase}
  \defaultfontfeatures[\rmfamily]{Ligatures=TeX,Scale=1}
\fi
\usepackage{lmodern}
\ifPDFTeX\else
  % xetex/luatex font selection
\fi
% Use upquote if available, for straight quotes in verbatim environments
\IfFileExists{upquote.sty}{\usepackage{upquote}}{}
\IfFileExists{microtype.sty}{% use microtype if available
  \usepackage[]{microtype}
  \UseMicrotypeSet[protrusion]{basicmath} % disable protrusion for tt fonts
}{}
\makeatletter
\@ifundefined{KOMAClassName}{% if non-KOMA class
  \IfFileExists{parskip.sty}{%
    \usepackage{parskip}
  }{% else
    \setlength{\parindent}{0pt}
    \setlength{\parskip}{6pt plus 2pt minus 1pt}}
}{% if KOMA class
  \KOMAoptions{parskip=half}}
\makeatother
\usepackage{color}
\usepackage{fancyvrb}
\newcommand{\VerbBar}{|}
\newcommand{\VERB}{\Verb[commandchars=\\\{\}]}
\DefineVerbatimEnvironment{Highlighting}{Verbatim}{commandchars=\\\{\}}
% Add ',fontsize=\small' for more characters per line
\usepackage{framed}
\definecolor{shadecolor}{RGB}{248,248,248}
\newenvironment{Shaded}{\begin{snugshade}}{\end{snugshade}}
\newcommand{\AlertTok}[1]{\textcolor[rgb]{0.94,0.16,0.16}{#1}}
\newcommand{\AnnotationTok}[1]{\textcolor[rgb]{0.56,0.35,0.01}{\textbf{\textit{#1}}}}
\newcommand{\AttributeTok}[1]{\textcolor[rgb]{0.13,0.29,0.53}{#1}}
\newcommand{\BaseNTok}[1]{\textcolor[rgb]{0.00,0.00,0.81}{#1}}
\newcommand{\BuiltInTok}[1]{#1}
\newcommand{\CharTok}[1]{\textcolor[rgb]{0.31,0.60,0.02}{#1}}
\newcommand{\CommentTok}[1]{\textcolor[rgb]{0.56,0.35,0.01}{\textit{#1}}}
\newcommand{\CommentVarTok}[1]{\textcolor[rgb]{0.56,0.35,0.01}{\textbf{\textit{#1}}}}
\newcommand{\ConstantTok}[1]{\textcolor[rgb]{0.56,0.35,0.01}{#1}}
\newcommand{\ControlFlowTok}[1]{\textcolor[rgb]{0.13,0.29,0.53}{\textbf{#1}}}
\newcommand{\DataTypeTok}[1]{\textcolor[rgb]{0.13,0.29,0.53}{#1}}
\newcommand{\DecValTok}[1]{\textcolor[rgb]{0.00,0.00,0.81}{#1}}
\newcommand{\DocumentationTok}[1]{\textcolor[rgb]{0.56,0.35,0.01}{\textbf{\textit{#1}}}}
\newcommand{\ErrorTok}[1]{\textcolor[rgb]{0.64,0.00,0.00}{\textbf{#1}}}
\newcommand{\ExtensionTok}[1]{#1}
\newcommand{\FloatTok}[1]{\textcolor[rgb]{0.00,0.00,0.81}{#1}}
\newcommand{\FunctionTok}[1]{\textcolor[rgb]{0.13,0.29,0.53}{\textbf{#1}}}
\newcommand{\ImportTok}[1]{#1}
\newcommand{\InformationTok}[1]{\textcolor[rgb]{0.56,0.35,0.01}{\textbf{\textit{#1}}}}
\newcommand{\KeywordTok}[1]{\textcolor[rgb]{0.13,0.29,0.53}{\textbf{#1}}}
\newcommand{\NormalTok}[1]{#1}
\newcommand{\OperatorTok}[1]{\textcolor[rgb]{0.81,0.36,0.00}{\textbf{#1}}}
\newcommand{\OtherTok}[1]{\textcolor[rgb]{0.56,0.35,0.01}{#1}}
\newcommand{\PreprocessorTok}[1]{\textcolor[rgb]{0.56,0.35,0.01}{\textit{#1}}}
\newcommand{\RegionMarkerTok}[1]{#1}
\newcommand{\SpecialCharTok}[1]{\textcolor[rgb]{0.81,0.36,0.00}{\textbf{#1}}}
\newcommand{\SpecialStringTok}[1]{\textcolor[rgb]{0.31,0.60,0.02}{#1}}
\newcommand{\StringTok}[1]{\textcolor[rgb]{0.31,0.60,0.02}{#1}}
\newcommand{\VariableTok}[1]{\textcolor[rgb]{0.00,0.00,0.00}{#1}}
\newcommand{\VerbatimStringTok}[1]{\textcolor[rgb]{0.31,0.60,0.02}{#1}}
\newcommand{\WarningTok}[1]{\textcolor[rgb]{0.56,0.35,0.01}{\textbf{\textit{#1}}}}
\usepackage{graphicx}
\makeatletter
\newsavebox\pandoc@box
\newcommand*\pandocbounded[1]{% scales image to fit in text height/width
  \sbox\pandoc@box{#1}%
  \Gscale@div\@tempa{\textheight}{\dimexpr\ht\pandoc@box+\dp\pandoc@box\relax}%
  \Gscale@div\@tempb{\linewidth}{\wd\pandoc@box}%
  \ifdim\@tempb\p@<\@tempa\p@\let\@tempa\@tempb\fi% select the smaller of both
  \ifdim\@tempa\p@<\p@\scalebox{\@tempa}{\usebox\pandoc@box}%
  \else\usebox{\pandoc@box}%
  \fi%
}
% Set default figure placement to htbp
\def\fps@figure{htbp}
\makeatother
\setlength{\emergencystretch}{3em} % prevent overfull lines
\providecommand{\tightlist}{%
  \setlength{\itemsep}{0pt}\setlength{\parskip}{0pt}}
\usepackage{bookmark}
\IfFileExists{xurl.sty}{\usepackage{xurl}}{} % add URL line breaks if available
\urlstyle{same}
\hypersetup{
  pdftitle={Likelihood Ratios},
  hidelinks,
  pdfcreator={LaTeX via pandoc}}

\title{Likelihood Ratios}
\author{}
\date{\vspace{-2.5em}2026-04-17}

\begin{document}
\maketitle

\subsection{Introduction}\label{introduction}

Likelihood ratios (LR+ and LR-) summarise the information content of a
diagnostic test in a way that is independent of disease prevalence.
Sensitivity and specificity describe how the test performs in
known-disease and known-no-disease populations, but they do not directly
tell a clinician what to believe after observing a positive or negative
test result; predictive values do, but they depend on prevalence and are
therefore not portable across populations. Likelihood ratios bridge this
gap: they combine multiplicatively with the pre-test odds to yield the
post-test odds via Bayes' theorem, making them the natural building
blocks of clinical Bayesian reasoning. Modern evidence-based-medicine
guidance and core teaching texts present LRs as the preferred
diagnostic-performance summary precisely because of this
prevalence-independence.

\subsection{Prerequisites}\label{prerequisites}

A working understanding of sensitivity and specificity, the relationship
between probability and odds, and Bayes' theorem in odds form.

\subsection{Theory}\label{theory}

The two basic likelihood ratios are

\[\mathrm{LR}^+ = \frac{\mathrm{Sens}}{1 - \mathrm{Spec}}, \qquad \mathrm{LR}^- = \frac{1 - \mathrm{Sens}}{\mathrm{Spec}}.\]

The Bayes-theorem update is

\[\mathrm{Odds}_{\text{post}} = \mathrm{Odds}_{\text{pre}} \times \mathrm{LR}.\]

Conventional clinical interpretation: \(\mathrm{LR}^+ > 10\) or
\(\mathrm{LR}^- < 0.1\) is often decisive; 5--10 or 0.1--0.2 is moderate
evidence; 2--5 or 0.2--0.5 is weak; near 1 is uninformative. Fagan's
nomogram graphically converts pre-test probability and LR directly to
post-test probability and is a useful bedside tool.

\subsection{Assumptions}\label{assumptions}

The same assumptions as for sensitivity and specificity: a reliable
gold-standard reference for disease status, accurate test
classification, and that the test characteristics estimated in one
population generalise to the patients to whom the LR is being applied.
Verification bias and spectrum bias both threaten this generalisation.

\subsection{R Implementation}\label{r-implementation}

\begin{Shaded}
\begin{Highlighting}[]
\FunctionTok{library}\NormalTok{(epiR)}

\NormalTok{tab }\OtherTok{\textless{}{-}} \FunctionTok{as.table}\NormalTok{(}\FunctionTok{matrix}\NormalTok{(}\FunctionTok{c}\NormalTok{(}\DecValTok{80}\NormalTok{, }\DecValTok{20}\NormalTok{,}
                         \DecValTok{20}\NormalTok{, }\DecValTok{880}\NormalTok{),}
                        \AttributeTok{nrow =} \DecValTok{2}\NormalTok{, }\AttributeTok{byrow =} \ConstantTok{FALSE}\NormalTok{,}
                        \AttributeTok{dimnames =} \FunctionTok{list}\NormalTok{(}\AttributeTok{Test =} \FunctionTok{c}\NormalTok{(}\StringTok{"+"}\NormalTok{, }\StringTok{"{-}"}\NormalTok{),}
                                        \AttributeTok{Disease =} \FunctionTok{c}\NormalTok{(}\StringTok{"yes"}\NormalTok{, }\StringTok{"no"}\NormalTok{))))}
\FunctionTok{epi.tests}\NormalTok{(tab)}\SpecialCharTok{$}\NormalTok{detail[}\FunctionTok{c}\NormalTok{(}\StringTok{"lrpos"}\NormalTok{, }\StringTok{"lrneg"}\NormalTok{), ]}

\NormalTok{sens }\OtherTok{\textless{}{-}} \FloatTok{0.8}\NormalTok{; spec }\OtherTok{\textless{}{-}} \FloatTok{0.978}
\NormalTok{lr\_pos }\OtherTok{\textless{}{-}}\NormalTok{ sens }\SpecialCharTok{/}\NormalTok{ (}\DecValTok{1} \SpecialCharTok{{-}}\NormalTok{ spec); lr\_neg }\OtherTok{\textless{}{-}}\NormalTok{ (}\DecValTok{1} \SpecialCharTok{{-}}\NormalTok{ sens) }\SpecialCharTok{/}\NormalTok{ spec}

\NormalTok{prior\_odds }\OtherTok{\textless{}{-}} \FloatTok{0.1} \SpecialCharTok{/} \FloatTok{0.9}
\NormalTok{post\_odds  }\OtherTok{\textless{}{-}}\NormalTok{ prior\_odds }\SpecialCharTok{*}\NormalTok{ lr\_pos}
\NormalTok{post\_prob  }\OtherTok{\textless{}{-}}\NormalTok{ post\_odds }\SpecialCharTok{/}\NormalTok{ (}\DecValTok{1} \SpecialCharTok{+}\NormalTok{ post\_odds)}
\FunctionTok{c}\NormalTok{(}\AttributeTok{lr\_pos =}\NormalTok{ lr\_pos, }\AttributeTok{lr\_neg =}\NormalTok{ lr\_neg, }\AttributeTok{post\_prob\_after\_pos =}\NormalTok{ post\_prob)}
\end{Highlighting}
\end{Shaded}

\subsection{Output \& Results}\label{output-results}

\texttt{epi.tests()} reports LR+ and LR- with their confidence intervals
from the input contingency table. The Bayesian update example shows how
a 10 \% pre-test probability rises to roughly 80 \% post-test after a
positive result with LR+ = 36, illustrating the multiplicative-on-odds
nature of the update.

\subsection{Interpretation}\label{interpretation}

A reporting sentence: ``The diagnostic test had sensitivity 80 \% (95 \%
CI 71 to 87 \%) and specificity 97.8 \% (96.4 to 98.7), corresponding to
LR+ = 36 (95 \% CI 16 to 81) and LR- = 0.20 (95 \% CI 0.13 to 0.31). A
positive test raised the pre-test probability of 10 \% to a post-test
probability of 80 \%, while a negative test reduced it to roughly 2 \%.
The test is therefore strongly informative in both directions for the
typical pre-test probability range of this clinical setting.'' Always
report LRs with CIs.

\subsection{Practical Tips}\label{practical-tips}

\begin{itemize}
\tightlist
\item
  Report likelihood ratios with their 95 \% confidence intervals; wide
  CIs indicate fragile diagnostic estimates and should temper
  interpretation, especially when small sample sizes or rare disease
  drive uncertainty in sensitivity or specificity.
\item
  For multi-category or ordinal tests (rating scales, semi-quantitative
  biomarker results), compute stratum-specific LRs for each score level
  rather than collapsing to a single binary LR; this preserves the
  information content of the gradation.
\item
  Likelihood ratios generalise across populations as long as the test
  characteristics (sensitivity, specificity) hold in the new population;
  this is their primary advantage over predictive values, which depend
  on local prevalence and do not transfer.
\item
  Clinical decision thresholds are often pre-specified in terms of
  required LR (e.g., LR+ ≥ 10 to justify initiating treatment, LR- ≤ 0.1
  to confidently rule out disease); building these thresholds into the
  diagnostic pathway is the operational analogue of Bayesian reasoning
  at the bedside.
\item
  Chain multiple test results by multiplying their LRs only if the tests
  are conditionally independent given disease status; in practice this
  assumption is often violated (a second test of the same type is
  correlated with the first), and joint LRs from a multivariable
  predictor are often more honest.
\item
  For complex multi-variable diagnostic tools (clinical prediction
  rules), the LR concept generalises naturally --- the rule's score
  corresponds to a stratum-specific LR --- and is a useful way to
  communicate the rule's discrimination at decision-relevant cut-points.
\end{itemize}

\subsection{R Packages Used}\label{r-packages-used}

\texttt{epiR::epi.tests()} for canonical sensitivity, specificity,
predictive values, and likelihood ratios with confidence intervals from
a 2 × 2 table; \texttt{epibasix} for an alternative interface;
\texttt{caret::confusionMatrix()} for ML-style classification metrics
including LRs; \texttt{pROC} for LRs across the full ROC operating
range; \texttt{bayesmeta} and related packages for Bayesian
meta-analytic pooling of likelihood ratios across diagnostic-accuracy
studies.

\subsection{For Reviewers}\label{for-reviewers}

\textbf{What to look for in a paper using this method.}

\begin{itemize}
\tightlist
\item
  \textbf{Common misapplications.}
\item
  \textbf{Diagnostics that should be reported but often aren't.}
\item
  \textbf{Red flags in tables and figures.}
\item
  \textbf{What to verify.}
\item
  \textbf{What an adequate Methods paragraph must contain.}
\end{itemize}

\subsection{See also --- labs in this
chapter}\label{see-also-labs-in-this-chapter}

\begin{itemize}
\tightlist
\item
  \href{labs/lab-diagnostic-testing-se-sp-ppv-npv-lr.Rmd}{Diagnostic
  testing: Se, Sp, PPV, NPV, LR}
\item
  \href{labs/lab-kappa-icc-bland-altman.Rmd}{Kappa, ICC, Bland--Altman}
\item
  \href{labs/lab-biomarker-statistics-youden-nri-decision-curves.Rmd}{Biomarker
  statistics (Youden, NRI, decision curves)}
\item
  \href{labs/lab-tripod-ai-fairness-auditing-reproducibility-at-scale.Rmd}{TRIPOD-AI,
  fairness auditing, reproducibility at scale}
\end{itemize}

\end{document}
